\documentclass[9pt]{elife}

\usepackage{amsmath,amssymb,mathtools}
\usepackage{siunitx}
\usepackage{graphicx}
\usepackage{booktabs}
\usepackage{hyperref}
\usepackage{tcolorbox}
\usepackage{bm}

% TITLE — candidate, in the "characterize validity" frame (replaces the old "exact/unbiased" titles).
% Alternative: "Where time-averaged likelihoods hold and where they break: macroscopic kinetic inference"
\title{Validity limits of time-averaged likelihoods for macroscopic kinetic inference}

\author[1*]{Luciano Moffatt}
\affil[1]{Instituto de Qu\'imica F\'isica de los Materiales, Medio Ambiente y Energ\'ia (INQUIMAE), CONICET; Facultad de Ciencias Exactas y Naturales, Universidad de Buenos Aires, Ciudad de Buenos Aires, C1428EHA, Argentina}
\corr{lmoffatt@qi.fcen.uba.ar}

\begin{document}
\maketitle

\begin{abstract}
Time-averaged measurements pervade biology, and the likelihoods used to interpret them are approximations whose reliability is rarely tested. We ask where such an approximation can be trusted and how it fails when it cannot. For macroscopic ion-channel currents, we simulate a two-state channel under a concentration jump and use the exact simulation as ground truth to judge a family of Gaussian likelihood approximations (non-recursive, recursive, interval-averaged, and their combinations). We measure faithfulness with likelihood-level diagnostics: standardized residuals, the score, and an information-distortion matrix that compares the covariance of the score with the model's own Fisher information. We then map how each approximation deforms inference as a function of the number of channels, the acquisition interval, and the instrumental noise. Only the interval-averaged recursive likelihood (MacroIR) stays calibrated across the practical regime; the others lose or misallocate information in ways that trace to the two Gaussian approximations they make, one on channel occupancy and one on the single-interval signal. We locate the multinomial, telegraphic, and Gaussian regimes where the approximations degrade, and we show the distortion carries a concrete consequence: it maps to a correction of the Bayesian evidence that ranks competing mechanisms. The result is a validity map for time-averaged likelihoods and a calibrated default, MacroIR, for macroscopic recordings.
\end{abstract}

% Impact statement (eLife requires one, <= 40 words):
% A validity map shows where time-averaged likelihood approximations for ion-channel currents can be trusted and where they distort inference, and identifies MacroIR as the one that stays calibrated across the practical regime.

% ================= SKELETON: titles only below. Fill-hints in % comments. =================

\section*{Introduction}
% The problem: measurements are interval-averaged; the likelihoods for them are approximations rarely tested.
% The question: not "is it exact" but "where does it hold and how does it fail".
% The asymmetry that licenses the approach: forward simulation is easy, the inverse (data -> likelihood) is hard, so the simulation is ground truth.
% Position: Comm Biol demonstration (Moffatt & Pierdominici-Sottile 2025) + Macro R lineage (Moffatt 2007) + Munch 2022 / Kalman prior art. This paper characterizes the likelihood-algorithm component.

\section*{The macroscopic interval likelihood and its approximations}
% Theory. Salvage the boundary-state derivation from archives/ (merged, revised3): interval average, E_2 kernel, Q(+)Q factorization, meta-state tcolorbox, the Kalman-like update and gain, and the emission-variance decomposition (eps^2/t + white + pink + N_ch[gating]); the E_3 variance kernel from revised3.
% Define the family NR / NMR / R / MR / IR and the two Gaussian approximations each makes (occupancy multinomial->Gaussian; single-interval signal->Gaussian). Salvage the full family derivation from archives/elife-macroir-merged.tex.

\section*{Diagnostics: testing a likelihood against its simulation}
% The four indicators: standardized residuals (mean 0, var 1, whiteness); score bias; information-distortion matrix C = H^{-1/2} J H^{-1/2} anchored on the Gaussian Fisher, decomposed into correlation vs sample/geometric distortion; direct empirical-vs-sandwich covariance test. Evaluate at the optimum. Salvage the score-test prose from archives/; the distortion matrix is new.

\section*{Results}

\subsection*{A minimal two-state model and the algorithm family}
% The clean 2-state C-O model, non-stationary pulse; why minimal (isolate the algorithm's own error).

\subsection*{Which approximations are calibrated: parameter recovery}
% Fig 2: MLE clouds + 3 ellipses (empirical / Gaussian-Fisher / sandwich). NR/NMR overconfident; IR calibrated. [regenerate on Gaussian rerun]

\subsection*{Time-resolved calibration and the score/Fisher checks}
% Fig 3: logL gap, R^2 residual, score bias, per-interval and accumulated J_t/F_t, score autocorrelation. [Gaussian-anchored]

\subsection*{Where the information lives: the Fisher falls to zero on relaxation}
% Fig 4 (key result): information about N_ch, k_on, i vanishes once the open count relaxes; k_off stays informative. [Gaussian-anchored]

\subsection*{A validity map across channels, interval, and noise}
% Fig 5: bias and distortion heatmaps over N_ch x K_off, all algorithms; the multinomial / telegraphic / Gaussian regimes. [regenerate on Gaussian rerun]

\subsection*{Decomposing the distortion: correlation versus sample}
% Fig 6: correlation vs sample/geometric distortion; where and why each approximation fails.

\section*{Discussion}
% MacroIR the calibrated default; the failure modes trace to the two approximations; ranking (R ~factor 2, MR strawman, NR/NMR inflation, NMR unbiased-but-underdispersed).
% MacroIR ~ time-augmented integrated Kalman filter (state as hypothesis, verified ~1e-8).
% The distortion is the correction to the Bayesian evidence (volume 1/2 log det C, effective-sample alpha* = p/tr C): motivation for the program's model-comparison component; derivation deferred.
% Research-program framing + explicit open doors: micro / >2 states / stationary / experimental data.

\section*{Materials and methods}

\subsection*{Minimal model and simulation}
\subsection*{The macroscopic likelihood approximations}
\subsection*{Diagnostics and the Gaussian-Fisher anchor}
\subsection*{Regime grid and reduction}
\subsection*{Reproducibility}
% Engine macro_dr @ pinned tag/hash; figure scripts; data handoff. See papers/macroir-elife-2025/09_carve_plan.md.

\section*{Data and code availability}
% GitHub repo (macroir-validity) + Zenodo DOI at submission; macro_dr engine referenced by pinned version.

\section*{Acknowledgements}

\section*{Author contributions}
% LM conceived the study, developed the theory, implemented the algorithms, analyzed data, and wrote the manuscript.

% \nocite{*}
\bibliography{biblio}

\end{document}
